Public cervical cytology benchmarks release many cell crops cut from far fewer microscope
fields or slide images, and are usually split at the crop level, so test cells share their
acquisition with training cells. A first study showed that this structure is pervasive and that
the shared parent is recognisable, but six per-architecture tests could not show that it inflates
accuracy. In a second, preregistered study (1,341 models) we test the question once, with a
mixed model sized by simulation: across twelve architectures and twenty new seeds, crop-level
splitting raised SIPaKMeD accuracy by \HoneEst\ percentage points (95\% CI \HoneBootLo--\HoneBootHi),
so parent-level splitting reveals \SipErrorRatio\ times the error rate. On CRIC, whose slide
images hold many cells of several classes, the inflation was \HtwoEst\ points. A factorial
that holds crop and field counts fixed shows that the inflation is leakage (\HthreeEst\ points)
and not the extra training diversity crop-level splits bring (diversity effect \DivEst\ points). Eight further registered
experiments (1,580 runs) locate the leaked signal by elimination. It needs no background and,
in an exploratory split, does not depend on shared pixels; it is not a slide-level label prior; and it is not carried by
linearly decodable field identity: erasing or adversarially suppressing that identity cut probe
recoverability by up to \RecoveryCutMax\% without moving the gap. It is instead tied to how a
field renders its cells: matching a test cell's colour statistics to an unseen field removed
most of it, although colour augmentation during training did not prevent it. Grad-CAM maps of
leaking and non-leaking models were indistinguishable. Neither probes, saliency maps, nor
augmentation can certify an evaluation; parent-disjoint splits can. We release the provenance
maps, splits, run records, and a risk calculator that predicts a benchmark's exposure from its
crop-to-parent table.
